\documentclass[aspectratio=169]{beamer}

% Configuración del tema y colores
\usetheme{Madrid}
\usecolortheme{default}

% Definición del color institucional UCB
\definecolor{ucbblue}{RGB}{0, 51, 102}

% Ajuste maestro de colores para Beamer (Soluciona el problema del título)
\setbeamercolor{structure}{fg=ucbblue}
\setbeamercolor*{palette primary}{fg=white,bg=ucbblue}
\setbeamercolor*{palette secondary}{fg=white,bg=ucbblue!80}
\setbeamercolor*{palette tertiary}{fg=white,bg=ucbblue!90}
\setbeamercolor{title}{fg=white, bg=ucbblue} % Texto blanco sobre fondo azul
\setbeamercolor{frametitle}{fg=white, bg=ucbblue}
\setbeamercolor{item}{fg=ucbblue}

% Paquetes estándar
\usepackage[utf8]{inputenc}
\usepackage[spanish]{babel}
\usepackage{amsmath, amssymb, bm}
\usepackage{tikz}
\usetikzlibrary{shapes.geometric, arrows.meta, positioning, calc, backgrounds}
\usepackage{listings}

% Configuración de código para Beamer
\lstset{
    language=Python,
    basicstyle=\ttfamily\scriptsize,
    keywordstyle=\color{blue}\bfseries,
    commentstyle=\color{green!50!black},
    stringstyle=\color{red},
    showstringspaces=false,
    frame=single,
    backgroundcolor=\color{gray!10},
    rulecolor=\color{gray!40},
}

% Información del documento
\title[Dualidad de SO(3)]{Definición Matemática de la Orientación en $\mathbb{R}^3$}
\subtitle{Robótica Industrial (IMT-342) — Universidad Católica Boliviana Tarija}
\author[B. Quiroga Turdera]{M.Sc. Ing. Bernardo Quiroga Turdera}
\date{Martes 11 de Agosto de 2026}

\begin{document}

\begin{frame}
    \titlepage
\end{frame}

% Diapositiva 1
\begin{frame}{Orientación Espacial y el Grupo $SO(3)$}
    \begin{columns}
        \begin{column}{0.65\textwidth}
            \textbf{Ecuación Central:}
            \begin{equation*}
                ^{A}R_B = \begin{bmatrix} ^{A}\hat{x}_B & ^{A}\hat{y}_B & ^{A}\hat{z}_B \end{bmatrix} = \begin{bmatrix} 
                \hat{x}_B \cdot \hat{x}_A & \hat{y}_B \cdot \hat{x}_A & \hat{z}_B \cdot \hat{x}_A \\ 
                \hat{x}_B \cdot \hat{y}_A & \hat{y}_B \cdot \hat{y}_A & \hat{z}_B \cdot \hat{y}_A \\ 
                \hat{x}_B \cdot \hat{z}_A & \hat{y}_B \cdot \hat{z}_A & \hat{z}_B \cdot \hat{z}_A 
                \end{bmatrix}
            \end{equation*}
            \vspace{0.2cm}
            
            \begin{block}{Concepto Clave}
                $^{A}R_B \in SO(3)$ es una matriz de $3 \times 3$ cuya columna $i$ representa la \textbf{proyección del eje unitario} del marco móvil $\{B\}$ sobre los ejes del marco fijo $\{A\}$.
            \end{block}
        \end{column}
        \begin{column}{0.35\textwidth}
            \centering
            \begin{tikzpicture}[scale=1.5,
                axis/.style={thick, ->, >=stealth},
                axis A/.style={axis, black},
                axis B/.style={axis, ucbblue}
            ]
                \draw[axis A] (0,0,0) -- (1.5,0,0) node[anchor=north]{$\hat{x}_A$};
                \draw[axis A] (0,0,0) -- (0,1.5,0) node[anchor=south]{$\hat{y}_A$};
                \draw[axis A] (0,0,0) -- (0,0,1.5) node[anchor=south]{$\hat{z}_A$};
                \node at (-0.75,-0.1) {$\{A\}$};
                
                \begin{scope}[rotate around z=30, rotate around y=20]
                    \draw[axis B] (0,0,0) -- (1.5,0,0) node[anchor=north]{$\hat{x}_B$};
                    \draw[axis B] (0,0,0) -- (0,1.5,0) node[anchor=south]{$\hat{y}_B$};
                    \draw[axis B] (0,0,0) -- (0,0,1.5) node[anchor=west]{$\hat{z}_B$};
                    \node[text=ucbblue] at (1.7,0.3,0.8) {$\{B\}$};
                \end{scope}
            \end{tikzpicture}
        \end{column}
    \end{columns}
\end{frame}

% Diapositiva 2
\begin{frame}{Las Dos Interpretaciones de $R \in SO(3)$}
    \begin{columns}[t]
        \begin{column}{0.48\textwidth}
            \begin{exampleblock}{\centering 1. Transformación Activa\\ (Operador Vectorial)}
                \begin{center}
                    \Large $p' = R \cdot p$
                \end{center}
                \vspace{0.2cm}
                \small Mueve un vector $p$ físico a una nueva posición $p'$ dentro del \textbf{mismo} sistema de referencia inercial $\{A\}$.
            \end{exampleblock}
        \end{column}
        
        \begin{column}{0.48\textwidth}
            \begin{block}{\centering 2. Transformación Pasiva\\ (Cambio de Marco)}
                \begin{center}
                    \Large $p_A = {}^{A}R_B \cdot p_B$
                \end{center}
                \vspace{0.2cm}
                \small Expresa un punto estático $p$ (con coordenadas locales $p_B$) en las coordenadas del sistema de referencia base $\{A\}$.
            \end{block}
        \end{column}
    \end{columns}
    
    \vspace{0.5cm}
    \begin{center}
        \textit{\textbf{Pregunta de Reflexión:} "Si un motor del robot gira $30^\circ$, ¿estamos aplicando una rotación activa al eslabón o una transformación pasiva al sistema de coordenadas de la herramienta?"}
    \end{center}
\end{frame}

% Diapositiva 3
\begin{frame}{Rotaciones Elementales sobre los Ejes Principales}
    \textbf{Ecuaciones Canónicas:}
    \vspace{0.1cm}
    \small
    \begin{equation*}
        R_x(\theta) = \begin{bmatrix} 1 & 0 & 0 \\ 0 & \cos\theta & -\sin\theta \\ 0 & \sin\theta & \cos\theta \end{bmatrix}, \quad
        R_y(\theta) = \begin{bmatrix} \cos\theta & 0 & \sin\theta \\ 0 & 1 & 0 \\ -\sin\theta & 0 & \cos\theta \end{bmatrix}, \quad
        R_z(\theta) = \begin{bmatrix} \cos\theta & -\sin\theta & 0 \\ \sin\theta & \cos\theta & 0 \\ 0 & 0 & 1 \end{bmatrix}
    \end{equation*}
    
    \vspace{0.3cm}
    \normalsize
    \textbf{Demostración de Ortogonalidad:}
    \begin{equation*}
        R^T R = I \implies R^{-1} = R^T \quad \text{y} \quad \det(R) = +1
    \end{equation*}
    
    \vspace{0.3cm}
    \begin{alertblock}{Propiedad Isométrica}
        La matriz conserva la longitud de los vectores (no deforma el espacio): \\ 
        \centerline{$\left\| R \cdot p \right\| = \left\| p \right\|$}
    \end{alertblock}
\end{frame}

% Diapositiva 4
\begin{frame}{Reglas de Multiplicación: Ejes Fijos vs. Ejes Móviles}
    \begin{center}
    \begin{tikzpicture}[
        node distance=1.5cm and 1cm,
        box/.style={rectangle, draw=ucbblue, thick, fill=ucbblue!10, text width=4cm, align=center, rounded corners, minimum height=1cm},
        root/.style={rectangle, draw=black, thick, fill=gray!20, text width=6cm, align=center, minimum height=0.8cm, font=\bfseries},
        arrow/.style={-{Stealth}, thick}
    ]
        % Nodos
        \node[root] (q) {¿Respecto a qué eje se realiza el giro?};
        
        \node[box, below left=1.2cm and -0.5cm of q] (fijo) {\textbf{Ejes Fijos} \\ de la Base $\{A\}$};
        \node[box, below right=1.2cm and -0.5cm of q] (movil) {\textbf{Ejes Móviles} \\ del Cuerpo $\{B\}$};
        
        \node[align=center, below=0.6cm of fijo] (pre) {\textbf{Pre-multiplicación}\\(IZQUIERDA)\\$R_{\text{total}} = R_{\text{nueva}} \cdot R_{\text{anterior}}$};
        \node[align=center, below=0.6cm of movil] (post) {\textbf{Post-multiplicación}\\(DERECHA)\\$R_{\text{total}} = R_{\text{anterior}} \cdot R_{\text{nueva}}$};
        
        % Conectores
        \draw[arrow] (q.south) -- (fijo.north);
        \draw[arrow] (q.south) -- (movil.north);
        \draw[arrow] (fijo.south) -- (pre.north);
        \draw[arrow] (movil.south) -- (post.north);
    \end{tikzpicture}
    \end{center}
    
    \vspace{0.2cm}
    \small \textbf{Justificación:} Moverse respecto a ejes fijos opera sobre la base inercial; moverse respecto a ejes locales encadena la cinemática de los eslabones anteriores.
\end{frame}

% Diapositiva 5
\begin{frame}[fragile]{Implementación Numérica y Verificación Isométrica}
\begin{lstlisting}
import numpy as np

# Angulo de 45 grados en radianes
theta = np.radians(45)
R_z = np.array([
    [np.cos(theta), -np.sin(theta), 0],
    [np.sin(theta),  np.cos(theta), 0],
    [0,              0,             1]
])

# Vector original p en el marco B
p_B = np.array([1.0, 2.0, 3.0])

# 1. Transformacion Pasiva: Cambio de base de B a A
p_A = np.dot(R_z, p_B)

# 2. Verificacion de Isometria (Preservacion de Norma)
print("Norma (p_B):", np.linalg.norm(p_B))
print("Norma transformada (p_A):", np.linalg.norm(p_A))
print("Ortogonal? (R^T*R==I):", np.allclose(np.dot(R_z.T, R_z), np.eye(3)))
\end{lstlisting}
\end{frame}

\end{document}